% ---------------------------------------------------------------------
% Methodology — first draft (STEPS.md §1.7).
% Self-contained section; input from main.tex via % ---------------------------------------------------------------------
% Methodology — first draft (STEPS.md §1.7).
% Self-contained section; input from main.tex via % ---------------------------------------------------------------------
% Methodology — first draft (STEPS.md §1.7).
% Self-contained section; input from main.tex via % ---------------------------------------------------------------------
% Methodology — first draft (STEPS.md §1.7).
% Self-contained section; input from main.tex via \input{methodology}.
% Appendix with the 6 reference prompts lives in appendix_prompts.tex.
% ---------------------------------------------------------------------

\section{Methodology}
\label{sec:methodology}

This section describes the unified pipeline used to run
$3\ \text{models} \times 3\ \text{domains} \times 2\ \text{prompting conditions} = 18$
experimental cells, and the mechanical steps we took to keep every
cell comparable. The key design commitment is that the only thing that
varies between two cells is the axis being studied --- the model, the
domain, or the prompting condition. Everything else (prompt shell,
generation configuration, iteration policy, validator) is controlled
by a single source of truth and verified on disk before and after the
SLURM runs.

\subsection{Pipeline overview}
\label{sec:methodology:pipeline}

Figure~\ref{fig:pipeline} shows the end-to-end flow. A parametric
SLURM entry point
(\texttt{scripts/slurm/run.sh}) accepts three environment variables
--- \texttt{MODEL}, \texttt{DOMAIN}, \texttt{CONDITION} --- and is submitted
18 times by \texttt{scripts/slurm/submit\_all.sh}. Each job loads one
model once, iterates over the 20 instances of the assigned domain,
and for every instance runs a prompt $\rightarrow$ generate $\rightarrow$
validate loop with up to 4 iterations under a stop-on-success policy
(Section~\ref{sec:methodology:iteration}). Per-iteration rows are
appended to \texttt{run\_metrics.csv}; the final plan is written to
\texttt{instance-XX\_plan.txt} with a metadata block that records
\texttt{first\_valid\_iter}.

\begin{figure}[t]
\centering
\begin{tikzpicture}[
    node distance=6mm and 10mm,
    box/.style={draw, rounded corners=2pt, align=center,
                minimum height=8mm, inner sep=4pt, font=\small},
    cell/.style={draw, rounded corners=2pt, align=center,
                 fill=black!3, minimum height=8mm, inner sep=4pt,
                 font=\small},
    arr/.style={-{Latex[length=2mm]}}
]
\node[box]                                      (cfg)   {\texttt{config.yml}\\(single source of truth)};
\node[box, below=of cfg]                        (sub)   {\texttt{submit\_all.sh}\\($3{\times}3{\times}2 = 18$ sbatches)};
\node[cell, below=of sub]                       (cell)  {SLURM cell --- \texttt{run.sh}\\$(M, D, C)$: 20 instances, up to 4 iters};
\node[box, below=of cell, xshift=-28mm]         (llm)   {LLM\\(vLLM, TP-split)};
\node[box, below=of cell]                       (val)   {VAL\\(\texttt{validate\_plan\_verbose})};
\node[box, below=of cell, xshift=28mm]          (met)   {Metrics\\(extractor $\to$ CSV)};
\node[box, below=of val]                        (out)   {\texttt{results/\{model\}/\{domain\}/\{condition\}/}\\($\star$ plans, $\star$ iter dumps, run\_metrics.csv)};

\draw[arr] (cfg)  -- (sub);
\draw[arr] (sub)  -- (cell);
\draw[arr] (cell) -- (llm);
\draw[arr] (cell) -- (val);
\draw[arr] (cell) -- (met);
\draw[arr] (llm)  |- ([yshift=-4mm]val.north) -- (val);
\draw[arr] (val)  -- (out);
\draw[arr] (met)  |- (out);
\end{tikzpicture}
\caption{Pipeline overview. \texttt{config.yml} is loaded once per job
and echoed to stdout at job start (Section~\ref{sec:methodology:genconfig}),
giving a post-hoc audit trail for the generation-configuration
uniformity claim.}
\label{fig:pipeline}
\end{figure}

\subsection{Model slate}

We evaluate three instruction-tuned models spanning roughly one order
of magnitude in parameter count and two families:

\begin{table}[h]
\centering
\small
\begin{tabular}{llll}
\toprule
Alias     & Repository                           & Family   & Profile \\
\midrule
llama8    & \texttt{meta-llama/Llama-3.1-8B-Instruct}  & Meta     & small (1$\times$ A100-80\,GB, 3\,h) \\
qwen25    & \texttt{Qwen/Qwen2.5-32B-Instruct}         & Alibaba  & mid   (2$\times$ A100-80\,GB TP=2, 10\,h) \\
llama33   & \texttt{meta-llama/Llama-3.3-70B-Instruct} & Meta     & large (2$\times$ A100-80\,GB TP=2, FP8, 14\,h) \\
\bottomrule
\end{tabular}
\caption{Model slate. \texttt{llama8} vs \texttt{llama33} pins the
intra-family scale axis; \texttt{qwen25} at 32B provides a cross-family
mid-scale reference point.}
\end{table}

The slate choice is motivated in the Experimental Design section; here
we only note that the SLURM profile associated with each model is
hard-coded in \texttt{run.sh} so that resubmissions cannot silently
drift onto a different GPU count, walltime, or quantization mode.

\subsection{Prompt harmonization}
\label{sec:methodology:prompts}

The two source projects used a single monolithic prompt file each,
branching internally on domain. We replaced this with a three-file
module under \texttt{src/prompts/}:

\begin{description}
\item[\texttt{shell.py}] Domain-agnostic primitives:
      \texttt{SYSTEM\_PROMPT\_PDDL} (the PDDL planner persona),
      \texttt{COT\_SCAFFOLD(description, problem)} (the chain-of-thought
      preamble), and \texttt{VALIDATION\_FEEDBACK\_TEMPLATE(val\_output,
      previous\_plan)} (the feedback used on iteration~2+).
\item[\texttt{descriptions.py}] One natural-language domain description per
      active domain, copied verbatim from the source projects to preserve
      prior phrasing. Dormant domains are declared as \texttt{None} and
      raise \texttt{NotImplementedError} at composition time.
\item[\texttt{compose.py}] The single entry point
      \texttt{build\_problem\_prompt(domain, condition, pddl\_domain, pddl\_problem)}.
      It takes \emph{no model argument}: byte-identity of prompts across
      models is therefore a structural property, not a runtime one.
\end{description}

Pseudocode for the composed prompt:

\begin{lstlisting}[basicstyle=\ttfamily\footnotesize]
SYSTEM_PROMPT_PDDL
<blank>
body     # = description  (baseline)  OR  COT_SCAFFOLD(description, problem)  (cot)
<blank>
=== DOMAIN DEFINITION ===
<pddl_domain>
<blank>
=== PROBLEM DEFINITION ===
<pddl_problem>
<blank>
Provide the plan now (one action per line).
\end{lstlisting}

The resulting prompt strings for \texttt{instance-01} under each
(domain, condition) pair are committed as fixtures at
\texttt{src/tests/fixtures/prompts/} and reproduced in
Appendix~\ref{app:prompts}.

\subsection{Prompting conditions}
\label{sec:methodology:conditions}

We study two conditions, differing only in the body slot of
Section~\ref{sec:methodology:prompts}:

\begin{itemize}
\item \textbf{baseline} --- the body is the bare domain description.
\item \textbf{cot} --- the body is the same description prepended by the
      CoT scaffold (\texttt{COT\_SCAFFOLD}), a 4-step instruction to parse
      initial state and goal, enumerate relevant actions, check
      preconditions, and minimise plan length, emitted inside
      \texttt{<think> \ldots </think>} tags.
\end{itemize}

All other tokens --- system prompt, domain PDDL, problem PDDL, trailing
instruction --- are byte-identical between the two conditions. The
\texttt{cot} prompt can be reconstructed exactly by inserting the
scaffold preamble into the \texttt{baseline} prompt, which is how
Section~\ref{sec:methodology:verification} defends the claim.

\subsection{Generation configuration uniformity}
\label{sec:methodology:genconfig}

Generation parameters are declared once, in \texttt{config.yml} under
the \texttt{generation} block, and reused by every cell:

\begin{lstlisting}[basicstyle=\ttfamily\footnotesize]
generation:
  temperature:    0.6
  top_p:          0.95
  top_k:          10
  max_new_tokens: 4096
  stop:           null
  seed:           42
\end{lstlisting}

No per-model overrides are permitted. At job start, \texttt{run.sh}
parses \texttt{config.yml} and echoes the resolved
\texttt{generation} and \texttt{iteration} blocks to stdout as a
single canonical JSON line tagged \texttt{"Generation config:"}. After
all 18 jobs complete, grepping the 18 stdout files and diffing these
lines gives a post-hoc, byte-exact audit that every cell ran with
identical sampling parameters (see Section~\ref{sec:methodology:verification}).

\subsection{Iteration policy --- stop-on-success}
\label{sec:methodology:iteration}

Each \texttt{(model, domain, condition, instance)} tuple is allowed up
to \texttt{max\_iterations = 4} attempts. The iteration loop is:

\begin{enumerate}
\item Build the problem prompt (first iteration) or the
      validation-feedback prompt (subsequent iterations).
\item Generate a plan from the LLM.
\item Write \texttt{instance-XX\_iter\_k.txt} to disk.
\item Run \texttt{validate\_plan\_verbose}; extract structured metrics.
\item Append one row to \texttt{run\_metrics.csv} with the columns
      \texttt{(model, domain, prompting\_condition, instance, iteration,
      plan\_text, plan\_len, solves\_problem, valid\_action\_percent,
      consecutive\_valid\_steps, logical\_violations, wallclock\_s,
      prompt\_tokens, completion\_tokens)}.
\item If the plan is valid, record \texttt{first\_valid\_iter = k} and
      break out of the loop.
\end{enumerate}

If all 4 iterations fail, the final \texttt{instance-XX\_plan.txt} is
written using the iter-4 plan with
\texttt{first\_valid\_iter = null} in the metadata block. This policy
is deliberate: it saves compute at the cost of a
well-understood selection bias --- iteration~$k$ is only populated by
instances that failed iterations $1..k{-}1$, so per-iteration metrics
across cells are not directly comparable without conditioning on
``reached iter $k$''. We flag this bias explicitly in every figure
that reports a per-iteration metric (Limitations).

\subsection{Validator setup}
\label{sec:methodology:validator}

Plan validation uses the VAL toolkit (Jan Dolejsi's PDDL-extension
build) pinned to a single commit SHA across local and HPC
environments. A thin helper in \texttt{src/utils/validator.py},

\begin{lstlisting}[basicstyle=\ttfamily\footnotesize]
validate_plan_verbose(domain_path, problem_path, plan_text)
    -> tuple[bool, str]
\end{lstlisting}

invokes \texttt{Validate -v} and returns \texttt{(is\_valid, raw\_stdout)}.
\texttt{src/utils/metrics\_extractor.py} (ported from Project A's
\texttt{utils/dataset\_generator.py}) parses the raw stdout into the
row schema above. A single validator plus a single parser, called from
a single place in the iteration loop, give us one source of truth for
every metric that appears in the Results section.

\subsection{Verification procedures}
\label{sec:methodology:verification}

Three checks run outside the SLURM critical path to defend the
uniformity claims made above:

\begin{enumerate}
\item \textbf{Prompt cross-model identity.}
      \texttt{scripts/verify/dump\_prompts.py} builds the prompt for
      \texttt{instance-01} under each \texttt{(model, domain, condition)}
      triple, hashes each output with SHA-256, and asserts that the three
      per-model hashes agree for every
      \texttt{(domain, condition)} combination. Because
      \texttt{build\_problem\_prompt} takes no model argument, the check
      is structurally guaranteed to pass; committing the SHA-256 digests
      and the resulting fixtures makes that guarantee auditable.
\item \textbf{Baseline-vs-CoT scaffold-only diff.}
      The same script reconstructs the \texttt{baseline} prompt by
      deleting the CoT scaffold block from the \texttt{cot} prompt and
      asserts the reconstruction is byte-equal to the committed
      \texttt{baseline} fixture. This rules out accidental drift of any
      shared boilerplate between the two conditions.
\item \textbf{Generation-configuration identity.}
      On Day~2 (after all 18 jobs have produced at least a startup log)
      we grep the \texttt{"Generation config:"} line from every stdout
      file; all 18 must be byte-identical. Any mismatch halts analysis
      until the affected cells are rerun with the canonical
      configuration.
\end{enumerate}

Together these checks operationalise the methodology: the 18 cells are
comparable because the pipeline that produced them was demonstrably
--- not merely intended to be --- uniform outside the three axes
of variation.
.
% Appendix with the 6 reference prompts lives in appendix_prompts.tex.
% ---------------------------------------------------------------------

\section{Methodology}
\label{sec:methodology}

This section describes the unified pipeline used to run
$3\ \text{models} \times 3\ \text{domains} \times 2\ \text{prompting conditions} = 18$
experimental cells, and the mechanical steps we took to keep every
cell comparable. The key design commitment is that the only thing that
varies between two cells is the axis being studied --- the model, the
domain, or the prompting condition. Everything else (prompt shell,
generation configuration, iteration policy, validator) is controlled
by a single source of truth and verified on disk before and after the
SLURM runs.

\subsection{Pipeline overview}
\label{sec:methodology:pipeline}

Figure~\ref{fig:pipeline} shows the end-to-end flow. A parametric
SLURM entry point
(\texttt{scripts/slurm/run.sh}) accepts three environment variables
--- \texttt{MODEL}, \texttt{DOMAIN}, \texttt{CONDITION} --- and is submitted
18 times by \texttt{scripts/slurm/submit\_all.sh}. Each job loads one
model once, iterates over the 20 instances of the assigned domain,
and for every instance runs a prompt $\rightarrow$ generate $\rightarrow$
validate loop with up to 4 iterations under a stop-on-success policy
(Section~\ref{sec:methodology:iteration}). Per-iteration rows are
appended to \texttt{run\_metrics.csv}; the final plan is written to
\texttt{instance-XX\_plan.txt} with a metadata block that records
\texttt{first\_valid\_iter}.

\begin{figure}[t]
\centering
\begin{tikzpicture}[
    node distance=6mm and 10mm,
    box/.style={draw, rounded corners=2pt, align=center,
                minimum height=8mm, inner sep=4pt, font=\small},
    cell/.style={draw, rounded corners=2pt, align=center,
                 fill=black!3, minimum height=8mm, inner sep=4pt,
                 font=\small},
    arr/.style={-{Latex[length=2mm]}}
]
\node[box]                                      (cfg)   {\texttt{config.yml}\\(single source of truth)};
\node[box, below=of cfg]                        (sub)   {\texttt{submit\_all.sh}\\($3{\times}3{\times}2 = 18$ sbatches)};
\node[cell, below=of sub]                       (cell)  {SLURM cell --- \texttt{run.sh}\\$(M, D, C)$: 20 instances, up to 4 iters};
\node[box, below=of cell, xshift=-28mm]         (llm)   {LLM\\(vLLM, TP-split)};
\node[box, below=of cell]                       (val)   {VAL\\(\texttt{validate\_plan\_verbose})};
\node[box, below=of cell, xshift=28mm]          (met)   {Metrics\\(extractor $\to$ CSV)};
\node[box, below=of val]                        (out)   {\texttt{results/\{model\}/\{domain\}/\{condition\}/}\\($\star$ plans, $\star$ iter dumps, run\_metrics.csv)};

\draw[arr] (cfg)  -- (sub);
\draw[arr] (sub)  -- (cell);
\draw[arr] (cell) -- (llm);
\draw[arr] (cell) -- (val);
\draw[arr] (cell) -- (met);
\draw[arr] (llm)  |- ([yshift=-4mm]val.north) -- (val);
\draw[arr] (val)  -- (out);
\draw[arr] (met)  |- (out);
\end{tikzpicture}
\caption{Pipeline overview. \texttt{config.yml} is loaded once per job
and echoed to stdout at job start (Section~\ref{sec:methodology:genconfig}),
giving a post-hoc audit trail for the generation-configuration
uniformity claim.}
\label{fig:pipeline}
\end{figure}

\subsection{Model slate}

We evaluate three instruction-tuned models spanning roughly one order
of magnitude in parameter count and two families:

\begin{table}[h]
\centering
\small
\begin{tabular}{llll}
\toprule
Alias     & Repository                           & Family   & Profile \\
\midrule
llama8    & \texttt{meta-llama/Llama-3.1-8B-Instruct}  & Meta     & small (1$\times$ A100-80\,GB, 3\,h) \\
qwen25    & \texttt{Qwen/Qwen2.5-32B-Instruct}         & Alibaba  & mid   (2$\times$ A100-80\,GB TP=2, 10\,h) \\
llama33   & \texttt{meta-llama/Llama-3.3-70B-Instruct} & Meta     & large (2$\times$ A100-80\,GB TP=2, FP8, 14\,h) \\
\bottomrule
\end{tabular}
\caption{Model slate. \texttt{llama8} vs \texttt{llama33} pins the
intra-family scale axis; \texttt{qwen25} at 32B provides a cross-family
mid-scale reference point.}
\end{table}

The slate choice is motivated in the Experimental Design section; here
we only note that the SLURM profile associated with each model is
hard-coded in \texttt{run.sh} so that resubmissions cannot silently
drift onto a different GPU count, walltime, or quantization mode.

\subsection{Prompt harmonization}
\label{sec:methodology:prompts}

The two source projects used a single monolithic prompt file each,
branching internally on domain. We replaced this with a three-file
module under \texttt{src/prompts/}:

\begin{description}
\item[\texttt{shell.py}] Domain-agnostic primitives:
      \texttt{SYSTEM\_PROMPT\_PDDL} (the PDDL planner persona),
      \texttt{COT\_SCAFFOLD(description, problem)} (the chain-of-thought
      preamble), and \texttt{VALIDATION\_FEEDBACK\_TEMPLATE(val\_output,
      previous\_plan)} (the feedback used on iteration~2+).
\item[\texttt{descriptions.py}] One natural-language domain description per
      active domain, copied verbatim from the source projects to preserve
      prior phrasing. Dormant domains are declared as \texttt{None} and
      raise \texttt{NotImplementedError} at composition time.
\item[\texttt{compose.py}] The single entry point
      \texttt{build\_problem\_prompt(domain, condition, pddl\_domain, pddl\_problem)}.
      It takes \emph{no model argument}: byte-identity of prompts across
      models is therefore a structural property, not a runtime one.
\end{description}

Pseudocode for the composed prompt:

\begin{lstlisting}[basicstyle=\ttfamily\footnotesize]
SYSTEM_PROMPT_PDDL
<blank>
body     # = description  (baseline)  OR  COT_SCAFFOLD(description, problem)  (cot)
<blank>
=== DOMAIN DEFINITION ===
<pddl_domain>
<blank>
=== PROBLEM DEFINITION ===
<pddl_problem>
<blank>
Provide the plan now (one action per line).
\end{lstlisting}

The resulting prompt strings for \texttt{instance-01} under each
(domain, condition) pair are committed as fixtures at
\texttt{src/tests/fixtures/prompts/} and reproduced in
Appendix~\ref{app:prompts}.

\subsection{Prompting conditions}
\label{sec:methodology:conditions}

We study two conditions, differing only in the body slot of
Section~\ref{sec:methodology:prompts}:

\begin{itemize}
\item \textbf{baseline} --- the body is the bare domain description.
\item \textbf{cot} --- the body is the same description prepended by the
      CoT scaffold (\texttt{COT\_SCAFFOLD}), a 4-step instruction to parse
      initial state and goal, enumerate relevant actions, check
      preconditions, and minimise plan length, emitted inside
      \texttt{<think> \ldots </think>} tags.
\end{itemize}

All other tokens --- system prompt, domain PDDL, problem PDDL, trailing
instruction --- are byte-identical between the two conditions. The
\texttt{cot} prompt can be reconstructed exactly by inserting the
scaffold preamble into the \texttt{baseline} prompt, which is how
Section~\ref{sec:methodology:verification} defends the claim.

\subsection{Generation configuration uniformity}
\label{sec:methodology:genconfig}

Generation parameters are declared once, in \texttt{config.yml} under
the \texttt{generation} block, and reused by every cell:

\begin{lstlisting}[basicstyle=\ttfamily\footnotesize]
generation:
  temperature:    0.6
  top_p:          0.95
  top_k:          10
  max_new_tokens: 4096
  stop:           null
  seed:           42
\end{lstlisting}

No per-model overrides are permitted. At job start, \texttt{run.sh}
parses \texttt{config.yml} and echoes the resolved
\texttt{generation} and \texttt{iteration} blocks to stdout as a
single canonical JSON line tagged \texttt{"Generation config:"}. After
all 18 jobs complete, grepping the 18 stdout files and diffing these
lines gives a post-hoc, byte-exact audit that every cell ran with
identical sampling parameters (see Section~\ref{sec:methodology:verification}).

\subsection{Iteration policy --- stop-on-success}
\label{sec:methodology:iteration}

Each \texttt{(model, domain, condition, instance)} tuple is allowed up
to \texttt{max\_iterations = 4} attempts. The iteration loop is:

\begin{enumerate}
\item Build the problem prompt (first iteration) or the
      validation-feedback prompt (subsequent iterations).
\item Generate a plan from the LLM.
\item Write \texttt{instance-XX\_iter\_k.txt} to disk.
\item Run \texttt{validate\_plan\_verbose}; extract structured metrics.
\item Append one row to \texttt{run\_metrics.csv} with the columns
      \texttt{(model, domain, prompting\_condition, instance, iteration,
      plan\_text, plan\_len, solves\_problem, valid\_action\_percent,
      consecutive\_valid\_steps, logical\_violations, wallclock\_s,
      prompt\_tokens, completion\_tokens)}.
\item If the plan is valid, record \texttt{first\_valid\_iter = k} and
      break out of the loop.
\end{enumerate}

If all 4 iterations fail, the final \texttt{instance-XX\_plan.txt} is
written using the iter-4 plan with
\texttt{first\_valid\_iter = null} in the metadata block. This policy
is deliberate: it saves compute at the cost of a
well-understood selection bias --- iteration~$k$ is only populated by
instances that failed iterations $1..k{-}1$, so per-iteration metrics
across cells are not directly comparable without conditioning on
``reached iter $k$''. We flag this bias explicitly in every figure
that reports a per-iteration metric (Limitations).

\subsection{Validator setup}
\label{sec:methodology:validator}

Plan validation uses the VAL toolkit (Jan Dolejsi's PDDL-extension
build) pinned to a single commit SHA across local and HPC
environments. A thin helper in \texttt{src/utils/validator.py},

\begin{lstlisting}[basicstyle=\ttfamily\footnotesize]
validate_plan_verbose(domain_path, problem_path, plan_text)
    -> tuple[bool, str]
\end{lstlisting}

invokes \texttt{Validate -v} and returns \texttt{(is\_valid, raw\_stdout)}.
\texttt{src/utils/metrics\_extractor.py} (ported from Project A's
\texttt{utils/dataset\_generator.py}) parses the raw stdout into the
row schema above. A single validator plus a single parser, called from
a single place in the iteration loop, give us one source of truth for
every metric that appears in the Results section.

\subsection{Verification procedures}
\label{sec:methodology:verification}

Three checks run outside the SLURM critical path to defend the
uniformity claims made above:

\begin{enumerate}
\item \textbf{Prompt cross-model identity.}
      \texttt{scripts/verify/dump\_prompts.py} builds the prompt for
      \texttt{instance-01} under each \texttt{(model, domain, condition)}
      triple, hashes each output with SHA-256, and asserts that the three
      per-model hashes agree for every
      \texttt{(domain, condition)} combination. Because
      \texttt{build\_problem\_prompt} takes no model argument, the check
      is structurally guaranteed to pass; committing the SHA-256 digests
      and the resulting fixtures makes that guarantee auditable.
\item \textbf{Baseline-vs-CoT scaffold-only diff.}
      The same script reconstructs the \texttt{baseline} prompt by
      deleting the CoT scaffold block from the \texttt{cot} prompt and
      asserts the reconstruction is byte-equal to the committed
      \texttt{baseline} fixture. This rules out accidental drift of any
      shared boilerplate between the two conditions.
\item \textbf{Generation-configuration identity.}
      On Day~2 (after all 18 jobs have produced at least a startup log)
      we grep the \texttt{"Generation config:"} line from every stdout
      file; all 18 must be byte-identical. Any mismatch halts analysis
      until the affected cells are rerun with the canonical
      configuration.
\end{enumerate}

Together these checks operationalise the methodology: the 18 cells are
comparable because the pipeline that produced them was demonstrably
--- not merely intended to be --- uniform outside the three axes
of variation.
.
% Appendix with the 6 reference prompts lives in appendix_prompts.tex.
% ---------------------------------------------------------------------

\section{Methodology}
\label{sec:methodology}

This section describes the unified pipeline used to run
$3\ \text{models} \times 3\ \text{domains} \times 2\ \text{prompting conditions} = 18$
experimental cells, and the mechanical steps we took to keep every
cell comparable. The key design commitment is that the only thing that
varies between two cells is the axis being studied --- the model, the
domain, or the prompting condition. Everything else (prompt shell,
generation configuration, iteration policy, validator) is controlled
by a single source of truth and verified on disk before and after the
SLURM runs.

\subsection{Pipeline overview}
\label{sec:methodology:pipeline}

Figure~\ref{fig:pipeline} shows the end-to-end flow. A parametric
SLURM entry point
(\texttt{scripts/slurm/run.sh}) accepts three environment variables
--- \texttt{MODEL}, \texttt{DOMAIN}, \texttt{CONDITION} --- and is submitted
18 times by \texttt{scripts/slurm/submit\_all.sh}. Each job loads one
model once, iterates over the 20 instances of the assigned domain,
and for every instance runs a prompt $\rightarrow$ generate $\rightarrow$
validate loop with up to 4 iterations under a stop-on-success policy
(Section~\ref{sec:methodology:iteration}). Per-iteration rows are
appended to \texttt{run\_metrics.csv}; the final plan is written to
\texttt{instance-XX\_plan.txt} with a metadata block that records
\texttt{first\_valid\_iter}.

\begin{figure}[t]
\centering
\begin{tikzpicture}[
    node distance=6mm and 10mm,
    box/.style={draw, rounded corners=2pt, align=center,
                minimum height=8mm, inner sep=4pt, font=\small},
    cell/.style={draw, rounded corners=2pt, align=center,
                 fill=black!3, minimum height=8mm, inner sep=4pt,
                 font=\small},
    arr/.style={-{Latex[length=2mm]}}
]
\node[box]                                      (cfg)   {\texttt{config.yml}\\(single source of truth)};
\node[box, below=of cfg]                        (sub)   {\texttt{submit\_all.sh}\\($3{\times}3{\times}2 = 18$ sbatches)};
\node[cell, below=of sub]                       (cell)  {SLURM cell --- \texttt{run.sh}\\$(M, D, C)$: 20 instances, up to 4 iters};
\node[box, below=of cell, xshift=-28mm]         (llm)   {LLM\\(vLLM, TP-split)};
\node[box, below=of cell]                       (val)   {VAL\\(\texttt{validate\_plan\_verbose})};
\node[box, below=of cell, xshift=28mm]          (met)   {Metrics\\(extractor $\to$ CSV)};
\node[box, below=of val]                        (out)   {\texttt{results/\{model\}/\{domain\}/\{condition\}/}\\($\star$ plans, $\star$ iter dumps, run\_metrics.csv)};

\draw[arr] (cfg)  -- (sub);
\draw[arr] (sub)  -- (cell);
\draw[arr] (cell) -- (llm);
\draw[arr] (cell) -- (val);
\draw[arr] (cell) -- (met);
\draw[arr] (llm)  |- ([yshift=-4mm]val.north) -- (val);
\draw[arr] (val)  -- (out);
\draw[arr] (met)  |- (out);
\end{tikzpicture}
\caption{Pipeline overview. \texttt{config.yml} is loaded once per job
and echoed to stdout at job start (Section~\ref{sec:methodology:genconfig}),
giving a post-hoc audit trail for the generation-configuration
uniformity claim.}
\label{fig:pipeline}
\end{figure}

\subsection{Model slate}

We evaluate three instruction-tuned models spanning roughly one order
of magnitude in parameter count and two families:

\begin{table}[h]
\centering
\small
\begin{tabular}{llll}
\toprule
Alias     & Repository                           & Family   & Profile \\
\midrule
llama8    & \texttt{meta-llama/Llama-3.1-8B-Instruct}  & Meta     & small (1$\times$ A100-80\,GB, 3\,h) \\
qwen25    & \texttt{Qwen/Qwen2.5-32B-Instruct}         & Alibaba  & mid   (2$\times$ A100-80\,GB TP=2, 10\,h) \\
llama33   & \texttt{meta-llama/Llama-3.3-70B-Instruct} & Meta     & large (2$\times$ A100-80\,GB TP=2, FP8, 14\,h) \\
\bottomrule
\end{tabular}
\caption{Model slate. \texttt{llama8} vs \texttt{llama33} pins the
intra-family scale axis; \texttt{qwen25} at 32B provides a cross-family
mid-scale reference point.}
\end{table}

The slate choice is motivated in the Experimental Design section; here
we only note that the SLURM profile associated with each model is
hard-coded in \texttt{run.sh} so that resubmissions cannot silently
drift onto a different GPU count, walltime, or quantization mode.

\subsection{Prompt harmonization}
\label{sec:methodology:prompts}

The two source projects used a single monolithic prompt file each,
branching internally on domain. We replaced this with a three-file
module under \texttt{src/prompts/}:

\begin{description}
\item[\texttt{shell.py}] Domain-agnostic primitives:
      \texttt{SYSTEM\_PROMPT\_PDDL} (the PDDL planner persona),
      \texttt{COT\_SCAFFOLD(description, problem)} (the chain-of-thought
      preamble), and \texttt{VALIDATION\_FEEDBACK\_TEMPLATE(val\_output,
      previous\_plan)} (the feedback used on iteration~2+).
\item[\texttt{descriptions.py}] One natural-language domain description per
      active domain, copied verbatim from the source projects to preserve
      prior phrasing. Dormant domains are declared as \texttt{None} and
      raise \texttt{NotImplementedError} at composition time.
\item[\texttt{compose.py}] The single entry point
      \texttt{build\_problem\_prompt(domain, condition, pddl\_domain, pddl\_problem)}.
      It takes \emph{no model argument}: byte-identity of prompts across
      models is therefore a structural property, not a runtime one.
\end{description}

Pseudocode for the composed prompt:

\begin{lstlisting}[basicstyle=\ttfamily\footnotesize]
SYSTEM_PROMPT_PDDL
<blank>
body     # = description  (baseline)  OR  COT_SCAFFOLD(description, problem)  (cot)
<blank>
=== DOMAIN DEFINITION ===
<pddl_domain>
<blank>
=== PROBLEM DEFINITION ===
<pddl_problem>
<blank>
Provide the plan now (one action per line).
\end{lstlisting}

The resulting prompt strings for \texttt{instance-01} under each
(domain, condition) pair are committed as fixtures at
\texttt{src/tests/fixtures/prompts/} and reproduced in
Appendix~\ref{app:prompts}.

\subsection{Prompting conditions}
\label{sec:methodology:conditions}

We study two conditions, differing only in the body slot of
Section~\ref{sec:methodology:prompts}:

\begin{itemize}
\item \textbf{baseline} --- the body is the bare domain description.
\item \textbf{cot} --- the body is the same description prepended by the
      CoT scaffold (\texttt{COT\_SCAFFOLD}), a 4-step instruction to parse
      initial state and goal, enumerate relevant actions, check
      preconditions, and minimise plan length, emitted inside
      \texttt{<think> \ldots </think>} tags.
\end{itemize}

All other tokens --- system prompt, domain PDDL, problem PDDL, trailing
instruction --- are byte-identical between the two conditions. The
\texttt{cot} prompt can be reconstructed exactly by inserting the
scaffold preamble into the \texttt{baseline} prompt, which is how
Section~\ref{sec:methodology:verification} defends the claim.

\subsection{Generation configuration uniformity}
\label{sec:methodology:genconfig}

Generation parameters are declared once, in \texttt{config.yml} under
the \texttt{generation} block, and reused by every cell:

\begin{lstlisting}[basicstyle=\ttfamily\footnotesize]
generation:
  temperature:    0.6
  top_p:          0.95
  top_k:          10
  max_new_tokens: 4096
  stop:           null
  seed:           42
\end{lstlisting}

No per-model overrides are permitted. At job start, \texttt{run.sh}
parses \texttt{config.yml} and echoes the resolved
\texttt{generation} and \texttt{iteration} blocks to stdout as a
single canonical JSON line tagged \texttt{"Generation config:"}. After
all 18 jobs complete, grepping the 18 stdout files and diffing these
lines gives a post-hoc, byte-exact audit that every cell ran with
identical sampling parameters (see Section~\ref{sec:methodology:verification}).

\subsection{Iteration policy --- stop-on-success}
\label{sec:methodology:iteration}

Each \texttt{(model, domain, condition, instance)} tuple is allowed up
to \texttt{max\_iterations = 4} attempts. The iteration loop is:

\begin{enumerate}
\item Build the problem prompt (first iteration) or the
      validation-feedback prompt (subsequent iterations).
\item Generate a plan from the LLM.
\item Write \texttt{instance-XX\_iter\_k.txt} to disk.
\item Run \texttt{validate\_plan\_verbose}; extract structured metrics.
\item Append one row to \texttt{run\_metrics.csv} with the columns
      \texttt{(model, domain, prompting\_condition, instance, iteration,
      plan\_text, plan\_len, solves\_problem, valid\_action\_percent,
      consecutive\_valid\_steps, logical\_violations, wallclock\_s,
      prompt\_tokens, completion\_tokens)}.
\item If the plan is valid, record \texttt{first\_valid\_iter = k} and
      break out of the loop.
\end{enumerate}

If all 4 iterations fail, the final \texttt{instance-XX\_plan.txt} is
written using the iter-4 plan with
\texttt{first\_valid\_iter = null} in the metadata block. This policy
is deliberate: it saves compute at the cost of a
well-understood selection bias --- iteration~$k$ is only populated by
instances that failed iterations $1..k{-}1$, so per-iteration metrics
across cells are not directly comparable without conditioning on
``reached iter $k$''. We flag this bias explicitly in every figure
that reports a per-iteration metric (Limitations).

\subsection{Validator setup}
\label{sec:methodology:validator}

Plan validation uses the VAL toolkit (Jan Dolejsi's PDDL-extension
build) pinned to a single commit SHA across local and HPC
environments. A thin helper in \texttt{src/utils/validator.py},

\begin{lstlisting}[basicstyle=\ttfamily\footnotesize]
validate_plan_verbose(domain_path, problem_path, plan_text)
    -> tuple[bool, str]
\end{lstlisting}

invokes \texttt{Validate -v} and returns \texttt{(is\_valid, raw\_stdout)}.
\texttt{src/utils/metrics\_extractor.py} (ported from Project A's
\texttt{utils/dataset\_generator.py}) parses the raw stdout into the
row schema above. A single validator plus a single parser, called from
a single place in the iteration loop, give us one source of truth for
every metric that appears in the Results section.

\subsection{Verification procedures}
\label{sec:methodology:verification}

Three checks run outside the SLURM critical path to defend the
uniformity claims made above:

\begin{enumerate}
\item \textbf{Prompt cross-model identity.}
      \texttt{scripts/verify/dump\_prompts.py} builds the prompt for
      \texttt{instance-01} under each \texttt{(model, domain, condition)}
      triple, hashes each output with SHA-256, and asserts that the three
      per-model hashes agree for every
      \texttt{(domain, condition)} combination. Because
      \texttt{build\_problem\_prompt} takes no model argument, the check
      is structurally guaranteed to pass; committing the SHA-256 digests
      and the resulting fixtures makes that guarantee auditable.
\item \textbf{Baseline-vs-CoT scaffold-only diff.}
      The same script reconstructs the \texttt{baseline} prompt by
      deleting the CoT scaffold block from the \texttt{cot} prompt and
      asserts the reconstruction is byte-equal to the committed
      \texttt{baseline} fixture. This rules out accidental drift of any
      shared boilerplate between the two conditions.
\item \textbf{Generation-configuration identity.}
      On Day~2 (after all 18 jobs have produced at least a startup log)
      we grep the \texttt{"Generation config:"} line from every stdout
      file; all 18 must be byte-identical. Any mismatch halts analysis
      until the affected cells are rerun with the canonical
      configuration.
\end{enumerate}

Together these checks operationalise the methodology: the 18 cells are
comparable because the pipeline that produced them was demonstrably
--- not merely intended to be --- uniform outside the three axes
of variation.
.
% Appendix with the 6 reference prompts lives in appendix_prompts.tex.
% ---------------------------------------------------------------------

\section{Methodology}
\label{sec:methodology}

This section describes the unified pipeline used to run
$3\ \text{models} \times 3\ \text{domains} \times 2\ \text{prompting conditions} = 18$
experimental cells, and the mechanical steps we took to keep every
cell comparable. The key design commitment is that the only thing that
varies between two cells is the axis being studied --- the model, the
domain, or the prompting condition. Everything else (prompt shell,
generation configuration, iteration policy, validator) is controlled
by a single source of truth and verified on disk before and after the
SLURM runs.

\subsection{Pipeline overview}
\label{sec:methodology:pipeline}

Figure~\ref{fig:pipeline} shows the end-to-end flow. A parametric
SLURM entry point
(\texttt{scripts/slurm/run.sh}) accepts three environment variables
--- \texttt{MODEL}, \texttt{DOMAIN}, \texttt{CONDITION} --- and is submitted
18 times by \texttt{scripts/slurm/submit\_all.sh}. Each job loads one
model once, iterates over the 20 instances of the assigned domain,
and for every instance runs a prompt $\rightarrow$ generate $\rightarrow$
validate loop with up to 4 iterations under a stop-on-success policy
(Section~\ref{sec:methodology:iteration}). Per-iteration rows are
appended to \texttt{run\_metrics.csv}; the final plan is written to
\texttt{instance-XX\_plan.txt} with a metadata block that records
\texttt{first\_valid\_iter}.

\begin{figure}[t]
\centering
\begin{tikzpicture}[
    node distance=6mm and 10mm,
    box/.style={draw, rounded corners=2pt, align=center,
                minimum height=8mm, inner sep=4pt, font=\small},
    cell/.style={draw, rounded corners=2pt, align=center,
                 fill=black!3, minimum height=8mm, inner sep=4pt,
                 font=\small},
    arr/.style={-{Latex[length=2mm]}}
]
\node[box]                                      (cfg)   {\texttt{config.yml}\\(single source of truth)};
\node[box, below=of cfg]                        (sub)   {\texttt{submit\_all.sh}\\($3{\times}3{\times}2 = 18$ sbatches)};
\node[cell, below=of sub]                       (cell)  {SLURM cell --- \texttt{run.sh}\\$(M, D, C)$: 20 instances, up to 4 iters};
\node[box, below=of cell, xshift=-28mm]         (llm)   {LLM\\(vLLM, TP-split)};
\node[box, below=of cell]                       (val)   {VAL\\(\texttt{validate\_plan\_verbose})};
\node[box, below=of cell, xshift=28mm]          (met)   {Metrics\\(extractor $\to$ CSV)};
\node[box, below=of val]                        (out)   {\texttt{results/\{model\}/\{domain\}/\{condition\}/}\\($\star$ plans, $\star$ iter dumps, run\_metrics.csv)};

\draw[arr] (cfg)  -- (sub);
\draw[arr] (sub)  -- (cell);
\draw[arr] (cell) -- (llm);
\draw[arr] (cell) -- (val);
\draw[arr] (cell) -- (met);
\draw[arr] (llm)  |- ([yshift=-4mm]val.north) -- (val);
\draw[arr] (val)  -- (out);
\draw[arr] (met)  |- (out);
\end{tikzpicture}
\caption{Pipeline overview. \texttt{config.yml} is loaded once per job
and echoed to stdout at job start (Section~\ref{sec:methodology:genconfig}),
giving a post-hoc audit trail for the generation-configuration
uniformity claim.}
\label{fig:pipeline}
\end{figure}

\subsection{Model slate}

We evaluate three instruction-tuned models spanning roughly one order
of magnitude in parameter count and two families:

\begin{table}[h]
\centering
\small
\begin{tabular}{llll}
\toprule
Alias     & Repository                           & Family   & Profile \\
\midrule
llama8    & \texttt{meta-llama/Llama-3.1-8B-Instruct}  & Meta     & small (1$\times$ A100-80\,GB, 3\,h) \\
qwen25    & \texttt{Qwen/Qwen2.5-32B-Instruct}         & Alibaba  & mid   (2$\times$ A100-80\,GB TP=2, 10\,h) \\
llama33   & \texttt{meta-llama/Llama-3.3-70B-Instruct} & Meta     & large (2$\times$ A100-80\,GB TP=2, FP8, 14\,h) \\
\bottomrule
\end{tabular}
\caption{Model slate. \texttt{llama8} vs \texttt{llama33} pins the
intra-family scale axis; \texttt{qwen25} at 32B provides a cross-family
mid-scale reference point.}
\end{table}

The slate choice is motivated in the Experimental Design section; here
we only note that the SLURM profile associated with each model is
hard-coded in \texttt{run.sh} so that resubmissions cannot silently
drift onto a different GPU count, walltime, or quantization mode.

\subsection{Prompt harmonization}
\label{sec:methodology:prompts}

The two source projects used a single monolithic prompt file each,
branching internally on domain. We replaced this with a three-file
module under \texttt{src/prompts/}:

\begin{description}
\item[\texttt{shell.py}] Domain-agnostic primitives:
      \texttt{SYSTEM\_PROMPT\_PDDL} (the PDDL planner persona),
      \texttt{COT\_SCAFFOLD(description, problem)} (the chain-of-thought
      preamble), and \texttt{VALIDATION\_FEEDBACK\_TEMPLATE(val\_output,
      previous\_plan)} (the feedback used on iteration~2+).
\item[\texttt{descriptions.py}] One natural-language domain description per
      active domain, copied verbatim from the source projects to preserve
      prior phrasing. Dormant domains are declared as \texttt{None} and
      raise \texttt{NotImplementedError} at composition time.
\item[\texttt{compose.py}] The single entry point
      \texttt{build\_problem\_prompt(domain, condition, pddl\_domain, pddl\_problem)}.
      It takes \emph{no model argument}: byte-identity of prompts across
      models is therefore a structural property, not a runtime one.
\end{description}

Pseudocode for the composed prompt:

\begin{lstlisting}[basicstyle=\ttfamily\footnotesize]
SYSTEM_PROMPT_PDDL
<blank>
body     # = description  (baseline)  OR  COT_SCAFFOLD(description, problem)  (cot)
<blank>
=== DOMAIN DEFINITION ===
<pddl_domain>
<blank>
=== PROBLEM DEFINITION ===
<pddl_problem>
<blank>
Provide the plan now (one action per line).
\end{lstlisting}

The resulting prompt strings for \texttt{instance-01} under each
(domain, condition) pair are committed as fixtures at
\texttt{src/tests/fixtures/prompts/} and reproduced in
Appendix~\ref{app:prompts}.

\subsection{Prompting conditions}
\label{sec:methodology:conditions}

We study two conditions, differing only in the body slot of
Section~\ref{sec:methodology:prompts}:

\begin{itemize}
\item \textbf{baseline} --- the body is the bare domain description.
\item \textbf{cot} --- the body is the same description prepended by the
      CoT scaffold (\texttt{COT\_SCAFFOLD}), a 4-step instruction to parse
      initial state and goal, enumerate relevant actions, check
      preconditions, and minimise plan length, emitted inside
      \texttt{<think> \ldots </think>} tags.
\end{itemize}

All other tokens --- system prompt, domain PDDL, problem PDDL, trailing
instruction --- are byte-identical between the two conditions. The
\texttt{cot} prompt can be reconstructed exactly by inserting the
scaffold preamble into the \texttt{baseline} prompt, which is how
Section~\ref{sec:methodology:verification} defends the claim.

\subsection{Generation configuration uniformity}
\label{sec:methodology:genconfig}

Generation parameters are declared once, in \texttt{config.yml} under
the \texttt{generation} block, and reused by every cell:

\begin{lstlisting}[basicstyle=\ttfamily\footnotesize]
generation:
  temperature:    0.6
  top_p:          0.95
  top_k:          10
  max_new_tokens: 4096
  stop:           null
  seed:           42
\end{lstlisting}

No per-model overrides are permitted. At job start, \texttt{run.sh}
parses \texttt{config.yml} and echoes the resolved
\texttt{generation} and \texttt{iteration} blocks to stdout as a
single canonical JSON line tagged \texttt{"Generation config:"}. After
all 18 jobs complete, grepping the 18 stdout files and diffing these
lines gives a post-hoc, byte-exact audit that every cell ran with
identical sampling parameters (see Section~\ref{sec:methodology:verification}).

\subsection{Iteration policy --- stop-on-success}
\label{sec:methodology:iteration}

Each \texttt{(model, domain, condition, instance)} tuple is allowed up
to \texttt{max\_iterations = 4} attempts. The iteration loop is:

\begin{enumerate}
\item Build the problem prompt (first iteration) or the
      validation-feedback prompt (subsequent iterations).
\item Generate a plan from the LLM.
\item Write \texttt{instance-XX\_iter\_k.txt} to disk.
\item Run \texttt{validate\_plan\_verbose}; extract structured metrics.
\item Append one row to \texttt{run\_metrics.csv} with the columns
      \texttt{(model, domain, prompting\_condition, instance, iteration,
      plan\_text, plan\_len, solves\_problem, valid\_action\_percent,
      consecutive\_valid\_steps, logical\_violations, wallclock\_s,
      prompt\_tokens, completion\_tokens)}.
\item If the plan is valid, record \texttt{first\_valid\_iter = k} and
      break out of the loop.
\end{enumerate}

If all 4 iterations fail, the final \texttt{instance-XX\_plan.txt} is
written using the iter-4 plan with
\texttt{first\_valid\_iter = null} in the metadata block. This policy
is deliberate: it saves compute at the cost of a
well-understood selection bias --- iteration~$k$ is only populated by
instances that failed iterations $1..k{-}1$, so per-iteration metrics
across cells are not directly comparable without conditioning on
``reached iter $k$''. We flag this bias explicitly in every figure
that reports a per-iteration metric (Limitations).

\subsection{Validator setup}
\label{sec:methodology:validator}

Plan validation uses the VAL toolkit (Jan Dolejsi's PDDL-extension
build) pinned to a single commit SHA across local and HPC
environments. A thin helper in \texttt{src/utils/validator.py},

\begin{lstlisting}[basicstyle=\ttfamily\footnotesize]
validate_plan_verbose(domain_path, problem_path, plan_text)
    -> tuple[bool, str]
\end{lstlisting}

invokes \texttt{Validate -v} and returns \texttt{(is\_valid, raw\_stdout)}.
\texttt{src/utils/metrics\_extractor.py} (ported from Project A's
\texttt{utils/dataset\_generator.py}) parses the raw stdout into the
row schema above. A single validator plus a single parser, called from
a single place in the iteration loop, give us one source of truth for
every metric that appears in the Results section.

\subsection{Verification procedures}
\label{sec:methodology:verification}

Three checks run outside the SLURM critical path to defend the
uniformity claims made above:

\begin{enumerate}
\item \textbf{Prompt cross-model identity.}
      \texttt{scripts/verify/dump\_prompts.py} builds the prompt for
      \texttt{instance-01} under each \texttt{(model, domain, condition)}
      triple, hashes each output with SHA-256, and asserts that the three
      per-model hashes agree for every
      \texttt{(domain, condition)} combination. Because
      \texttt{build\_problem\_prompt} takes no model argument, the check
      is structurally guaranteed to pass; committing the SHA-256 digests
      and the resulting fixtures makes that guarantee auditable.
\item \textbf{Baseline-vs-CoT scaffold-only diff.}
      The same script reconstructs the \texttt{baseline} prompt by
      deleting the CoT scaffold block from the \texttt{cot} prompt and
      asserts the reconstruction is byte-equal to the committed
      \texttt{baseline} fixture. This rules out accidental drift of any
      shared boilerplate between the two conditions.
\item \textbf{Generation-configuration identity.}
      On Day~2 (after all 18 jobs have produced at least a startup log)
      we grep the \texttt{"Generation config:"} line from every stdout
      file; all 18 must be byte-identical. Any mismatch halts analysis
      until the affected cells are rerun with the canonical
      configuration.
\end{enumerate}

Together these checks operationalise the methodology: the 18 cells are
comparable because the pipeline that produced them was demonstrably
--- not merely intended to be --- uniform outside the three axes
of variation.
